\documentclass[a4paper]{article}
\usepackage{times}
\usepackage[utf8]{inputenc}
\usepackage[T1]{fontenc}
\usepackage{selinput}
\usepackage{upquote}
\usepackage[margin=2cm, rmargin=4cm, tmargin=3cm]{geometry}
\usepackage{tcolorbox}
\usepackage{xspace}
\usepackage{amsmath}
\usepackage{float}
\usepackage[french]{babel}
\usepackage{hyperref}
\usepackage{xurl}
\usepackage{fontawesome6}
\usepackage{enumitem}
\usepackage{marginnote}
\usepackage{ulem}
\usepackage[yyyymmdd,hhmmss]{datetime}
\usepackage{ulem}
\usepackage{environ}
\usepackage{graphicx}
%\usepackage[top=Bcm, bottom=Hcm, outer=Ccm, inner=Acm, heightrounded, marginparwidth=Ecm, marginparsep=Dcm]{geometry}
\usepackage{courier}



\newtcolorbox{Example}[1]{colback=white,left=20pt,colframe=slideblue,fonttitle=\bfseries,title=#1}
\newtcolorbox{Solutions}[1]{colback=white,left=20pt,colframe=green,fonttitle=\bfseries,title=#1}
\newtcolorbox{Conseils}[1]{colback=white,left=20pt,colframe=slideblue,fonttitle=\bfseries,title=#1}
\newtcolorbox{Warning}[1]{colback=white,left=20pt,colframe=warning,fonttitle=\bfseries,title=#1}
\newtcolorbox{Note}[1]{colback=white,left=20pt,colframe=grey,fonttitle=\bfseries,title=#1}
\newtcolorbox{Dessin}[1]{colback=white,left=20pt,colframe=grey,fonttitle=\bfseries,title=#1}

\setlength\parindent{0pt}

  %Exercice environment
  \newcounter{exercice}
  \newenvironment{Exercice}[1][]
  {
  \par
  \stepcounter{exercice}\textbf{Question \arabic{exercice}:} (\faClock \enskip \textit{#1})
  }
  {\bigskip}
  

% Title
\newcommand{\titre}{\begin{center}
  \section*{Algorithmes et Pensée Computationnelle}
\end{center}}
\newcommand{\cours}[1]
{\begin{center} 
  \textit{#1}\\
\end{center}
  }



% \newcommand{\exemple}[1]{\newline~\textbf{Exemple :} #1}
\newcommand{\todo}[1]{\textcolor{red}{#1}}
%\newcommand{\attention}[1]{\newline\faExclamationTriangle~\textbf{Attention :} #1}

% Documentation url (escape \# in the TP document)
\newcommand{\documentation}[1]{\faBookOpen~Documentation : \href{#1}{#1}}

% Clef API
\newcommand{\apikey}[1]{\faKey~Clé API : \lstinline{#1}}
\newcommand{\apiendpoint}[1]{\faGlobe~Url de base de l'API \href{#1}{#1}}

%Listing Python style
\usepackage{color}
\definecolor{slideblue}{RGB}{33,131,189}
\definecolor{green}{RGB}{0,190,100}
\definecolor{blue}{RGB}{121,142,213}
\definecolor{grey}{RGB}{120,120,120}
\definecolor{warning}{RGB}{235,186,1}

\usepackage{listings}
\lstdefinelanguage{texte}{
    keywordstyle=\color{black},
    numbers=none,
    frame=none,
    literate=
           {é}{{\'e}}1
           {è}{{\`e}}1
           {ê}{{\^e}}1
           {à}{{\`a}}1
           {â}{{\^a}}1
           {ù}{{\`u}}1
           {ü}{{\"u}}1
           {î}{{\^i}}1
           {ï}{{\"i}}1
           {ë}{{\"e}}1
           {Ç}{{\,C}}1
           {ç}{{\,c}}1
           {ô}{{\^o}}1,
    columns=fullflexible,keepspaces,
	breaklines=true,
	breakatwhitespace=true,
}
\lstset{
  language=Python,
	%basicstyle=\bfseries\footnotesize,
  basicstyle=\footnotesize\ttfamily\small,
	breaklines=true,
	breakatwhitespace=true,
	commentstyle=\color{grey},
	stringstyle=\color{slideblue},
  keywordstyle=\color{slideblue},
	morekeywords={with, as, True, False, Fraction, Float, join, index, None, main, argparse, self, sort, __eq__, __add__, __mul__ ,__ne__, __radd__, __del__, __ge__, __gt__, split, os, endswith, is_file, scandir, @classmethod, this},
	deletekeywords={id},
	showspaces=false,
	showstringspaces=false,
	columns=fullflexible,keepspaces,
	literate=
           {é}{{\'e}}1
           {è}{{\`e}}1
           {ê}{{\^e}}1
           {à}{{\`a}}1
           {â}{{\^a}}1
           {ù}{{\`u}}1
           {ü}{{\"u}}1
           {î}{{\^i}}1
           {ï}{{\"i}}1
           {ë}{{\"e}}1
           {Ç}{{\,C}}1
           {ç}{{\,c}}1
           {ô}{{\^o}}1,
    numbers=left,
}

\newtcbox{\mybox}{nobeforeafter,colframe=white,colback=slideblue,boxrule=0.5pt,arc=1.5pt, boxsep=0pt,left=2pt,right=2pt,top=2pt,bottom=2pt,tcbox raise base}
\newcommand{\projet}{\mybox{\textcolor{white}{\small projet}}\xspace}
\newcommand{\optionnel}{\mybox{\textcolor{white}{\small Optionnel}}\xspace}
\newcommand{\advanced}{\mybox{\textcolor{white}{\small Pour aller plus loin}}\xspace}
\newcommand{\auto}{\mybox{\textcolor{white}{\small Auto-évaluation}}\xspace}


\usepackage{environ}
\newif\ifShowSolution
\NewEnviron{solution}{
  \ifShowSolution
	\begin{Solutions}{\faTerminal \enskip Solution}
		\BODY
	\end{Solutions}
  \fi}

\usepackage{environ}
\newif\ifShowExample
\NewEnviron{exemple}{
  \ifShowExample
	\begin{Example}{\faTerminal \enskip Exemple}
		\BODY
	\end{Example}
  \fi}

\usepackage{environ}
 \newif\ifShowConseil
  \NewEnviron{conseil}{
    \ifShowConseil
    \begin{Conseils}{\faLightbulb \quad Conseil}
      \BODY
    \end{Conseils}

    \fi}

    \usepackage{environ}
  \newif\ifShowWarning
  \NewEnviron{attention}{
    \ifShowWarning
    \begin{Warning}{\faTriangleExclamation \quad Attention}
      \BODY
    \end{Warning}

    \fi}

  \newif\ifShowNote
  \NewEnviron{note}[1]{
    \ifShowNote
    \begin{Note}{\faEye \quad Comment added by #1 at \date - \currenttime}
      \BODY
    \end{Note}

    \fi}  

    \newif\ifShowDessin
    \NewEnviron{dessin}[1]{
      \ifShowDessin
      \begin{Dessin}{\faPen \quad Figure}
        \BODY
      \end{Dessin}
  
      \fi}  

%\newcommand{\Conseil}[1]{\ifShowIndice\ \newline\faLightbulb[regular]~#1\fi}


\usepackage{listings}
\usepackage{array}
\newcolumntype{C}[1]{>{\centering\let\newline\\\arraybackslash\hspace{0pt}}m{#1}}

\begin{document}

% Change the following values to true to show the solutions or/and the hints
\ShowSolutionfalse
\ShowConseilfalse
\titre
\cours{Programmation de base - suite}

Le but de cette séance est d'aborder des notions de base en programmation et de consolider les connaissances acquises lors de la séance de TP précédente.
Au terme de cette séance, l'étudiant sera capable de:
\begin{itemize}
    \item manipuler des chaînes de caractères,
    \item utiliser des branchements conditionnels.
\end{itemize}


\section{Manipulation de chaînes de caractères}

\begin{Exercice}[5 minutes] \textbf{Manipulation des chaînes de caractères (indexation) (Java ou Python)}\\
   Créez une variable \texttt{mon\_mot} de type chaîne de caractères et affectez-lui la valeur \textbf{``Hard But Cool !!''}. À l'aide de l'indexation, créez ensuite une variable \texttt{premiere} contenant le premier caractère de \texttt{mon\_mot}, puis une variable \texttt{derniere} contenant son dernier caractère. Affichez les valeurs obtenues pour \texttt{premiere} et \texttt{derniere}. \\
   
    \begin{conseil}
      	Pour Python, utilisez \lstinline{[]}, et pour Java, utilisez la fonction \lstinline{substring()} ainsi que la fonction \lstinline{length()} qui permet d'obtenir la taille d'un élément.
        
    \end{conseil}
    \begin{solution}
    
    \textbf{Python:}
    \lstinputlisting{resources/solutions/question9.py}
    
    \textbf{\\Java: (à l'intérieur de la méthode main)}
    \lstinputlisting{resources/solutions/question9.java}
           
    \end{solution}   
\end{Exercice}

\begin{Exercice}[5 minutes] \textbf{Manipulation des chaînes de caractères (indexation 2) (Java ou Python)}\\
   Gardez votre variable, \texttt{mon\_mot}, et créez une variable \texttt{lettre\_5} contenant la cinquième lettre de \texttt{mon\_mot} en utilisant l'indexation. Créez ensuite une variable \texttt{lettre\_9\_13} contenant les lettres 9, 10, 11, 12, 13 de \texttt{mon\_mot}. Afficher les résultats et voyez ce que vous obtenez.  \\
   
    \begin{conseil}
      	Attention, ici les espaces comptent comme des lettres ! \\

		Pour Python, utilisez \lstinline{[:]}, et pour Java, utilisez la fonction \lstinline{substring()}.
        
    \end{conseil}
    \begin{solution}
    
    \textbf{Python:}    
    \lstinputlisting{resources/solutions/question10.py}
    
    \textbf{\\Java: (à l'intérieur de la méthode main)}
    \lstinputlisting{resources/solutions/question10.java}
           
    \end{solution}   
\end{Exercice}
\begin{Exercice}[10 minutes] \textbf{Manipulation des chaînes de caractères (Java ou Python) \optionnel}\\
   Il est possible d'obtenir la longueur d'une chaîne de caractères (ou d'une liste ou d'un dictionnaire) en utilisant la fonction \lstinline{len()}. Gardez votre variable \texttt{mon\_mot} et créez une nouvelle variable nommée \texttt{ln\_mon\_mot} contenant le nombre de caractères de la variable \texttt{mon\_mot}, puis une nouvelle variable \texttt{moitie} contenant la première moitié de la variable \texttt{mon\_mot} (utilisez la variable que vous venez de créer). Affichez le résultat.   \\
   
    \begin{conseil}
      	La fonction présentée dans l'énoncé de la question n'est valable que pour Python. L'équivalent pour Java est la fonction \lstinline{length()}.
        
    \end{conseil}
    \begin{solution}
    
    \textbf{Python}:
    \lstinputlisting{resources/solutions/question11.py}
    
    \textbf{\\Java: (à l'intérieur de la méthode main)}
    \lstinputlisting{resources/solutions/question11.java}
    \end{solution}   
\end{Exercice}

\newpage

\section{Conditions}
% TODO: Exercice à compléter

\begin{Exercice}[5 minutes] \textbf{Conditions (Python)}\\
    Qu'affiche le programme suivant? \\
    
    \lstinputlisting{resources/questions/Question12.py}

     \begin{conseil}
         En Python, la fonction \lstinline{not} renvoie l'opposé d'une valeur booléenne. Par exemple, \lstinline{not(False)} renverra \lstinline{True}. En Java, on utilise un point d'exclamation avant la valeur.
     \end{conseil}
     \begin{solution}
        Benoît: Professeur du cours d'Algorithmes et Pensée Computationnelle.
     \end{solution}   
 \end{Exercice}

 \begin{Exercice}[10 minutes] \textbf{Branchement conditionnel en Java}\\
   Qu'affiche le programme suivant? 
   
   \lstinputlisting{resources/solutions/question12.java}
    
     \begin{conseil}
          \begin{itemize}
             \item \lstinline{break} indique que l'on sort de l'accolade. Les cas suivants ne seront pas traités.
             \item L'absence de \lstinline{break} indique que l'on va rentrer dans tous les cas suivants, jusqu'à enfin atteindre un \lstinline{break}.
             \item Lorsque l'on pose \lstinline{case n} où \lstinline{n} est un nombre cela est équivalent au test \lstinline{n == numero_mois}. Ce test est aussi valable si on cherche à comparer des chaînes de caractères (par exemple si \lstinline{numero_mois = "Juin"}, à ce moment là \lstinline{n} sera aussi une chaîne de caractères).
          \end{itemize}
         
     \end{conseil}
     \begin{solution}
     
     Juillet
     
    Aout
    
    Décembre \\
    
    Explications:
    \begin{itemize}
             \item Comme le \lstinline{case 7} ne contient pas de \lstinline{break} et modifie \lstinline{numero_mois}, la lecture du code va continuer.
             \item On rentre dans le \lstinline{case 9}, qui contient un \lstinline{break}. Le \lstinline{numero_mois} sera aussi modifié mais cela ne sera pas important car on sort de l'accolade et les cas suivants ne seront pas traités.
       \end{itemize}
     \end{solution}   
 \end{Exercice}

\end{document}
